\section{Abordagens para Alocação de Fatias}

Diferentes estratégias podem ser utilizadas para realizar o fatiamento e adaptar a utilização dos recursos às condições apresentadas pela rede. Essas estratégias podem adotar diferentes critérios para definir as fatias, resultando em comportamentos distintos quanto à utilização dos recursos disponíveis. Nesse contexto, \citeonline{gomes2020strategies} apresentam quatro estratégias gulosas e determinísticas para realizar o fatiamento considerando as variações da demanda ao longo do dia: o \textit{Per Time Calculation} (PETIC), o \textit{Previous Slice Reuse} (PSR), o \textit{Slice Encourage} (SE) e o \textit{Comparing Slices Saturation} (CSS). Este trabalho toma o PETIC como estratégia de referência e o compara a três implementações de famílias de métodos distintas, que não constam entre as estratégias desse artigo: a meta-heurística populacional \textit{Ant Colony Optimization} (ACO), a meta-heurística de trajetória \textit{Simulated Annealing} (SA) e o aprendizado por reforço com \textit{Q-Learning}. As quatro abordagens são descritas a seguir.


\subsection{Per Time Calculation (PETIC)}

O Per Time Calculation (PETIC) é um algoritmo proposto para realizar a alocação de fatias considerando as variações na demanda de largura de banda dos clientes ao longo do tempo. Para cada período analisado, o algoritmo busca definir uma fatia capaz de atender aos requisitos do cliente, selecionando os recursos disponíveis na infraestrutura física. Dessa forma, seu processo de alocação considera o estado dos recursos da rede e a demanda apresentada em cada período \cite{gomes2018fatiamento}.

O PETIC é caracterizado como um algoritmo guloso e determinístico, realizando suas decisões de alocação a partir das condições apresentadas pela infraestrutura em cada momento. Para definir uma fatia, o algoritmo considera a largura de banda disponível nos enlaces e a largura de banda requisitada pelo cliente, atribuindo pesos aos recursos que podem atender à solicitação. A partir desses pesos, é realizado um processo de busca para encontrar os caminhos que formarão a fatia \cite{gomes2018fatiamento}.

Durante o processo de alocação, os enlaces que não possuem largura de banda suficiente para atender à requisição são desconsiderados. Para os enlaces restantes, o cálculo dos pesos considera a relação entre a largura de banda disponível e a largura de banda requisitada, favorecendo a utilização de recursos que apresentem maior capacidade disponível. Com isso, o algoritmo busca evitar a saturação dos enlaces e preservar recursos que possam ser utilizados em novas solicitações \cite{gomes2018fatiamento}.

Além da disponibilidade dos recursos no momento da alocação, o PETIC considera informações da fatia utilizada anteriormente. Caso a demanda do período corrente não supere a largura de banda alocada à fatia no período anterior, os enlaces dessa fatia recebem um peso muito próximo de zero, favorecendo sua reutilização. Essa característica busca diminuir alterações na estrutura da fatia entre diferentes períodos, evitando modificações desnecessárias quando a configuração anterior continua adequada às necessidades do cliente \cite{gomes2018fatiamento}.

Do ponto de vista operacional, o PETIC executa, para cada período, os seguintes passos: (i) atribui um peso a cada enlace, conforme a relação entre a banda disponível e a banda requisitada, ou o peso reduzido quando o enlace pertence à fatia anterior reaproveitável; (ii) descarta os enlaces cuja banda disponível é insuficiente para a demanda; e (iii) executa uma busca de menor caminho, no estilo do algoritmo de Dijkstra, sobre os enlaces restantes, retornando o caminho de menor custo total. Se não existir caminho, a requisição é considerada não atendida naquele período. O caminho encontrado é então armazenado e serve de base para a alocação do período seguinte \cite{gomes2018fatiamento}.

\subsection{Ant Colony Optimization (ACO)}

O algoritmo Ant Colony Optimization (ACO) é uma meta-heurística originada do Ant System, proposto por \citeonline{dorigo1996ant}, baseado no comportamento coletivo de formigas na busca por alimentos, principalmente no uso de informações por meio dos feromônios. Foi formulado como um método de otimização combinatória e inicialmente aplicado a problemas como o Traveling Salesman Problem (TSP). Posteriormente, diferentes algoritmos baseados no comportamento de colônias de formigas foram generalizados sob a estrutura da meta-heurística ACO \cite{dorigo1999ant}.

Fundamentalmente, o ACO funciona com base na construção incremental de soluções por um conjunto de formigas artificiais. Cada formiga seleciona probabilisticamente um dos componentes disponíveis para continuar a construção de uma solução. Essa escolha se baseia em duas informações principais: a trilha de feromônio $\tau$, que representa a experiência acumulada pelas formigas em decisões anteriores, e a informação heurística $\eta$, que indica a atratividade de uma escolha de acordo com as características do problema. A importância desses fatores é controlada pelos parâmetros $\alpha$ e $\beta$, sendo que as decisões com maiores valores de feromônio e de informação heurística têm maior tendência de serem selecionadas \cite{dorigo1999ant}. Após cada formiga realizar suas escolhas e finalizar a construção de sua solução, ocorre a atualização das trilhas de feromônio. Nesse processo, soluções de melhor qualidade podem receber um reforço na quantidade de feromônio, enquanto a evaporação reduz gradualmente os valores existentes, diminuindo a influência de decisões anteriores nas próximas iterações \cite{dorigo1999ant}.

No contexto da alocação de recursos e da virtualização de redes, abordagens baseadas em ACO têm sido investigadas como alternativas para problemas de otimização relacionados ao mapeamento e à alocação de recursos. \citeonline{fajjari2011vnACO} propuseram uma adaptação do ACO para o problema de Virtual Network Embedding (VNE), considerando a alocação de redes virtuais sobre uma infraestrutura física e avaliando a taxa de rejeição das requisições, o custo, a receita e a utilização média dos enlaces. Em outro contexto, \citeonline{wen2023power} utilizaram uma abordagem baseada em colônia de formigas para auxiliar na decisão de associação entre fluxos e slices, buscando maximizar o retorno de qualidade de serviço obtido por unidade de investimento em potência (ROI).

Considerando suas características, o ACO pode ser adaptado para solucionar problemas de fatiamento de redes tendo as possíveis alocações de recursos como soluções do processo de otimização. Nesse contexto, as formigas artificiais podem explorar diferentes configurações para atender às demandas existentes, enquanto os feromônios podem favorecer recursos presentes em soluções de melhor qualidade. Em ambientes de demanda elástica, esse processo pode considerar as alterações nas necessidades dos usuários ao longo do tempo, permitindo a busca por novas configurações de fatias conforme o estado dos recursos e as demandas apresentadas em cada período.
\subsection{Simulated Annealing (SA)}

O Simulated Annealing (SA) é uma técnica de otimização inspirada no processo físico de recozimento de materiais, onde um material é submetido a uma temperatura alta e posteriormente resfriado de forma gradual. Em temperaturas elevadas, a maior agitação térmica permite que os átomos se arranjem em configurações que temporariamente geram mais energia. Enquanto perde temperatura, essas transições tornam-se menos prováveis, estabilizando o sistema em configurações menos energéticas. Baseado nisso, \citeonline{kirkpatrick1983optimization} conseguiram criar uma analogia entre esse processo e problemas de otimização. 

Em processos de otimização, o SA funciona a partir da geração de soluções vizinhas com base em uma configuração atual. Quando uma solução apresenta um resultado melhor, ela é aceita como solução atual. Mas diferente de uma solução gulosa, que só aceita soluções melhores que a atual, o SA também permite, de forma probabilística, a aceitação temporária de soluções de maior custo. Essa característica deriva do procedimento de Metropolis, no qual uma alteração desfavorável pode ser aceita segundo uma probabilidade que diminui em função da magnitude da piora e da temperatura atual do sistema \cite{kirkpatrick1983optimization}.

De uma forma mais simples, considerando uma variação de custo \(\Delta\) e uma temperatura \(T\), a probabilidade de aceitação de uma solução de maior custo pode ser expressa por:
\[
P = e^{-\Delta/T}
\]
Assim, quando está mais quente, o algoritmo se mostra mais propenso a aceitar soluções temporariamente piores, fazendo com que seja possível explorar diferentes regiões do espaço em busca de resultados melhores. De acordo com a redução da temperatura, essa probabilidade diminui, deixando o processo progressivamente mais seletivo e direcionado ao refinamento das soluções encontradas. A possibilidade de aceitar temporariamente soluções de maior custo é uma das principais características do SA, pois reduz a tendência da busca permanecer restrita a ótimos locais. 

\citeonline{kirkpatrick1983optimization} demonstram a aplicabilidade da abordagem em diferentes problemas combinatórios, incluindo particionamento, posicionamento de componentes eletrônicos e roteamento das conexões entre eles, além do problema do caixeiro-viajante. \citeonline{de2021slice} empregam a técnica em um cenário de \textit{network slicing} 5G com o objetivo de reduzir o orçamento de potência da Rede de Acesso por Rádio, considerando requisitos distintos de capacidade, latência e confiabilidade das fatias. Esse emprego do SA em um problema de otimização de recursos associado ao network slicing evidencia a possibilidade de adaptação da meta-heurística a diferentes representações de solução e funções objetivo, desde que sejam definidos adequadamente os movimentos de vizinhança e as restrições do problema.





\subsection{Q-Learning} O Q-Learning é uma das abordagens mais difundidas de Aprendizado por Reforço (\textit{Reinforcement Learning} -- RL) livre de modelo (model-free), onde o conhecimento adquirido pelo agente é estruturado na chamada tabela Q (Q-Table). Essa estrutura de dados armazena os valores Q($s$,$a$), conhecidos como valores ação-estado (quality values). Cada linha da tabela representa um estado $s$ possível do ambiente, enquanto cada coluna corresponde a uma ação $a$ que o agente pode tomar. O valor armazenado na célula ($s$, $a$) mensura a utilidade ou a expectativa de recompensa futura acumulada ao escolher a ação $a$ estando no estado $s$ e, a partir do estado resultante, seguir uma política ótima.

Com uma probabilidade $\epsilon$, o algoritmo escolhe uma ação de forma aleatória. Isso permite que o agente conheça novas regiões do espaço de estados e descubra caminhos potencialmente melhores que ainda não foram testados. Para decidir entre testar uma ação inedita ou utilizar uma ação cujo retorno já é conhecido, o algoritmo implementa estratégias de balanço entre exploração (\textit{exploration}) e explotação (\textit{exploitation}). A abordagem mais comum é a política $\epsilon$-gulosa (\textit{$\epsilon$-greedy}). A regra de atualização apresentada a seguir foi proposta por \citeonline{watkins1992q}.

No paradigma de Aprendizagem por Reforço (\textit{Reinforcement Learning} -- RL), o processo de tomada de decisão é classicamente modelado pela interação contínua entre um agente inteligente e o ambiente dinâmico. A cada ciclo de decisão, a partir de um estado corrente $s \in \mathcal{S}$, o agente executa uma ação $a \in \mathcal{A}$, transita para um novo estado $s' \in \mathcal{S}$ e recebe como resposta um sinal escalar de reforço ou recompensa imediata $r \in \mathbb{R}$. 

Com base nessa experiência de transição, a função de ação-valor é refinada iterativamente na estrutura tabular de conhecimento (\textit{Q-table}) por meio da regra de atualização do algoritmo \textit{Q-learning}, expressa pela Equação~\eqref{eq:qlearning_geral}:

\begin{equation}
    Q(s, a) \leftarrow Q(s, a) + \alpha \left[ r + \gamma \max_{a' \in \mathcal{A}} Q(s', a') - Q(s, a) \right]
    \label{eq:qlearning_geral}
\end{equation}

\noindent Em que os termos e hiperparâmetros são definidos como:
\begin{itemize}
    \item $Q(s, a)$: estimativa do valor esperado do retorno cumulativo ao adotar a ação $a$ estando no estado $s$;
    \item $\alpha \in (0, 1]$: a \textbf{taxa de aprendizado} (\textit{learning rate}), responsável por ponderar o impacto das informações recém-adquiridas frente ao conhecimento acumulado historicamente;
    \item $r$: o reforço imediato obtido após a transição do estado $s$ para $s'$;
    \item $\gamma \in [0, 1]$: o \textbf{fator de desconto} (\textit{discount factor}), que calibra a relevância das recompensas futuras esperadas em relação ao benefício imediato;
    \item $\max_{a' \in \mathcal{A}} Q(s', a')$: a estimativa do valor ótimo da função de ação-valor avaliada a partir do estado subsequente $s'$.
\end{itemize}

No domínio de redes de computadores, o Q-Learning pode ser empregado para que o agente aprenda, por meio de tentativa e erro, as ações que proporcionam o melhor desempenho frente às dinâmicas observadas no ambiente. Ao modelar o problema de alocação de recursos como um Processo de Decisão de Markov (MDP), o estado $s$, o conjunto de ações $a$ e a função de recompensa $r$ precisam ser definidos de acordo com o problema. A literatura apresenta diferentes escolhas: por exemplo, \citeonline{kim2019reinforcement} utilizam o Q-Learning tabular para o ajuste dinâmico dos recursos adquiridos por um inquilino (\textit{tenant}) em um cenário de fatiamento com tráfego elástico e inelástico, adotando como estado a proporção de fluxos inelásticos, a demanda de banda, a razão entre os preços de compra e venda e a zona de tempo, e como recompensa o lucro do inquilino penalizado pela violação de requisitos de qualidade de serviço.

Neste trabalho, o problema de encontrar o caminho que compõe uma fatia é modelado de forma distinta, como descrito no Capítulo~\ref{chp:metodologia}: o estado é o nó em que o agente se encontra, cada ação corresponde à escolha de um nó vizinho alcançável por um enlace com banda suficiente, e a recompensa é o negativo do custo do enlace percorrido. Essa modelagem é própria deste trabalho, e não decorre das propostas citadas. O Q-Learning é, assim, utilizado como um buscador de caminhos treinado a cada alocação, em vez de um agente que acumula aprendizado entre períodos.

\section{Trabalhos Relacionados}

Pesquisas sobre elasticidade de recursos tratam camadas e problemas distintos. \citeonline{pasin2013infaestrutura} estudam o ajuste autonômico de réplicas de servidores em uma infraestrutura distribuída sujeita a variações de carga. \citeonline{maciel2020gerenciamento} propõe uma arquitetura de gerenciamento e orquestração de fatias elásticas que integra recursos de nuvem e conectividade SDN. Lima Filho~(\citeyear{limafilho2026controle}) avalia o compartilhamento de blocos de recursos em redes de acesso 5G Open RAN por meio de uma xApp. Esses trabalhos tratam, respectivamente, da elasticidade de servidores, da orquestração de fatias e dos recursos de rádio; nenhum deles trata a seleção de caminhos em topologias de transporte em malha considerada aqui.

No fatiamento elástico de redes, \citeonline{gomes2020strategies} analisam as estratégias PETIC, PSR, SE e CSS sob variações temporais da demanda. Este trabalho toma o PETIC como referência exata para o custo aditivo implementado e acrescenta ACO e Q-Learning como buscas aproximadas desse mesmo custo, além do SA com função objetivo que inclui equilíbrio de utilização e extensão de caminho. A contribuição é experimental: disponibiliza um simulador comum e reprodutível e estima, com cenários independentes e comparações pareadas, como essas implementações se comportam sob as mesmas requisições e restrições de capacidade. Não se propõe um novo algoritmo de alocação nem uma avaliação geral das meta-heurísticas ou do aprendizado por reforço.

